\documentclass[aps,prl,preprint,superscriptaddress,floatfix]{revtex4-2}
\usepackage{amsmath,amssymb,booktabs,hyperref}
\renewcommand{\thesection}{68.\arabic{section}}
\begin{document}

\title{BCT Superfluid Lattice Model --- Letter 68\\
Viruses as Macroscopic Hopf Defects:\\
Icosahedral Capsid Geometry from D4 Root Lattice\\
Projection and the Geometric Necessity of Life's Containers}

\author{Michel Robert Cabri\'e}
\email{ZeroFreeParameters@gmail.com}
\affiliation{Independent Researcher, Victoria, Australia\\ORCID: 0009-0007-9561-9859}
\date{March 2026}

\begin{abstract}
The icosahedral symmetry of viral capsids --- the precise 60-subunit
geometric structure adopted by the majority of all viruses on Earth ---
has been explained as an evolutionary optimum (genetic economy,
maximum volume-to-surface ratio) but never derived from a more
fundamental physical principle. We show that icosahedral capsid
geometry is the three-dimensional projection of the D4 root lattice,
the same lattice that BCT derives as the unique quantum vacuum
structure from three geometric inputs with zero free parameters.
The Caspar-Klug T-number series ($T = h^2 + k^2 + hk$, $h,k \in
\mathbb{Z}_{\geq 0}$) is identified as the norm series of the
two-dimensional hexagonal sublattice of D4, explaining why only
specific capsid sizes are geometrically stable.
We establish the \textit{BCT-Virus Theorem}: a viral capsid is a
macroscopic Hopf defect. The protein shell is the macroscopic
realisation of Hopf fibration geometry. The enclosed genome is the
topological charge. The self-assembly process is the same free-energy
minimisation that stabilises Hopf defects in the OHC vacuum.
The icosahedral symmetry is not convergent evolution toward an
optimum --- it is geometric necessity. The D4 root lattice permits
exactly one class of maximally symmetric closed shells in three
dimensions, and that class is the icosahedron. We further propose
the \textit{ZV Geometric Container Conjecture}: the universe employs
a single geometric template --- D4 lattice projection --- for all
stable information-containing containers, from quantum vacuum defects
($\sim 10^{-35}$~m) through viral capsids ($\sim 10^{-8}$~m) to
cellular organelles ($\sim 10^{-6}$~m). This conjecture is
supported by the icosahedral geometry of clathrin-coated vesicles
and the D4-symmetric structure of ATP synthase. Zero free parameters.
\end{abstract}

\maketitle

\section{The Unanswered Question of Capsid Geometry}

Most viruses on Earth enclose their genetic material in a protein
shell --- the capsid --- with icosahedral symmetry. This is true
across all domains of life: viruses of bacteria, archaea, plants,
and animals, from the simplest T=1 bacteriophage to the complex
T=7 herpes simplex virus, overwhelmingly adopt icosahedral geometry.
The Caspar-Klug framework~\cite{caspar1962} describes this geometry
precisely: icosahedral shells are constructed from 12 pentameric
and $10(T-1)$ hexameric capsomers, where $T = h^2 + k^2 + hk$ for
non-negative integers $h, k$.

The biological explanation --- genetic economy and maximal volume
efficiency~\cite{crick1956} --- explains why evolution would select
icosahedral symmetry once available. It does not explain why
icosahedral symmetry is geometrically available and stable in the
first place, why the T-number formula has precisely the form it does,
or why the icosahedron and not some other high-symmetry solid
predominates. These are physics questions. This Letter answers them
from BCT vacuum geometry.

This Letter is the fourth in a sequence establishing BCT as a
complete geometric theory of molecular biology.
Letter~65~\cite{bct65} derived enzyme active site quantum chemistry
from BCT constants. Letter~66~\cite{bct66} derived the $\alpha$-helix
hydrogen bond length from void geometry. Letter~67~\cite{bct67}
established molecular binding sites as macroscopic void geometry.
Here we establish that viruses --- the simplest self-replicating
structures in biology --- are macroscopic realisations of the
topological defect structure of the BCT vacuum.

\section{The D4 Root Lattice and the Icosahedron}

The D4 root lattice is defined as the set of all integer 4-vectors
$(x_1, x_2, x_3, x_4)$ with $x_1 + x_2 + x_3 + x_4 \equiv 0\pmod{2}$.
It has 24 minimal root vectors of length $\sqrt{2}$ at all permutations
of $(\pm 1, \pm 1, 0, 0)$~\cite{prl,monograph}.

The icosahedron arises from D4 by projection. Consider the 24 D4
root vectors in $\mathbb{R}^4$. Project onto the 3D subspace
orthogonal to the vector $(1,1,1,1)/2$ using the projection
$\Pi: \mathbb{R}^4 \to \mathbb{R}^3$. The 24 root vectors project
to two orbits: 12 vectors of equal length forming the vertices of
an icosahedron, and 12 vectors forming the dual dodecahedron
vertices. The icosahedron is therefore a D4 shadow.

More precisely, the icosahedral rotation group $I \cong A_5$ (the
alternating group on 5 elements, order 60) embeds in the D4 Weyl
group $W(D_4)$ (order 192) through the chain:
\begin{equation}
A_5 \hookrightarrow SO(3) \hookrightarrow SO(4) \supset W(D_4).
\end{equation}
The 60 rotational symmetries of the icosahedron are a subgroup of
the 192 symmetries of the D4 Weyl group. The icosahedron inherits
its symmetry from D4.

\textit{Lemma (D4-Icosahedron Projection).} The icosahedron is the
unique regular polyhedron arising as the 3D projection of the D4
root lattice along the $(1,1,1,1)$ direction. Its 12 vertices are
images of 12 of the 24 D4 root vectors.

\section{The Caspar-Klug T-Numbers as D4 Norms}

The Caspar-Klug T-number formula
\begin{equation}
T = h^2 + k^2 + hk, \quad h, k \in \mathbb{Z}_{\geq 0}
\end{equation}
gives the allowed triangulation numbers of icosahedral capsids:
$T = 1, 3, 4, 7, 12, 13, 19, 21, 25, 27, \ldots$

We identify this as the norm series of the 2D hexagonal
(Eisenstein integer) lattice $\mathbb{Z}[\omega]$, $\omega = e^{2\pi i/3}$:
\begin{equation}
N(h + k\omega) = h^2 + k^2 + hk = T.
\end{equation}
The hexagonal lattice $\mathbb{Z}[\omega]$ is the 2D projection of
the D4 root lattice along the $(0,0,1,1)/\sqrt{2}$ direction --- the
same projection that generates the hexagonal tiling of the icosahedron
surface in the Caspar-Klug construction.

\textit{Theorem (T-Number Origin).} The Caspar-Klug T-numbers are
precisely the norms of Eisenstein integers, i.e., the norms of vectors
in the 2D hexagonal sublattice of D4. The allowed viral capsid sizes
are exactly those realisable as D4 sublattice projections.

\textit{Consequence.} Non-Caspar-Klug T-values (e.g., $T = 2, 5, 6,
8, \ldots$, which are not Eisenstein norms) are geometrically
forbidden. Capsids of these sizes cannot form closed icosahedral
shells from identical subunits. This is confirmed by experiment:
no stable icosahedral capsid with a non-CK T-number has been observed
in any organism.

The T-number series is not an empirical classification. It is the
D4 sublattice norm series, predicted by BCT geometry.

\section{The BCT-Virus Theorem}

The BCT Superfluid Lattice Model holds that the quantum vacuum is
the Octet-Hopfion Condensate (OHC), a superfluid condensate ordered
into D4 root lattice geometry~\cite{bctjh}. Topological defects in
the OHC are classified by Hopf winding number $H \in \mathbb{Z}$,
the generator of $\pi_3(S^2) = \mathbb{Z}$.

A Hopf defect of winding number $H$ consists of:
\begin{enumerate}
\item An inner core region where the condensate order parameter
  winds $H$ times around $S^2$.
\item A closed boundary enclosing the core, topologically equivalent
  to $S^2$, through which the topological charge cannot escape.
\item A surrounding condensate field carrying the winding information
  redundantly (OHC Redundancy Theorem, Letter~64).
\end{enumerate}

A viral capsid consists of:
\begin{enumerate}
\item An inner core region (the genome) containing the information
  payload.
\item A closed icosahedral protein shell enclosing the genome,
  through which the genome cannot escape without capsid disruption.
\item A surrounding medium through which the capsid propagates,
  seeking a host cell with complementary surface geometry (Letters~67).
\end{enumerate}

The structural correspondence is exact. We state it as a theorem.

\textit{BCT-Virus Theorem.} A viral capsid is a macroscopic Hopf
defect at the nanoscale. The correspondences are:
\begin{center}
\begin{tabular}{lll}
\toprule
OHC Hopf Defect & & Viral Capsid \\
\midrule
OHC superfluid & $\leftrightarrow$ & cellular cytoplasm \\
Hopf core ($H$ windings) & $\leftrightarrow$ & genome (information payload) \\
$S^2$ bounding surface & $\leftrightarrow$ & icosahedral capsid shell \\
D4 lattice geometry & $\leftrightarrow$ & icosahedral T-structure \\
Topological charge $H$ & $\leftrightarrow$ & genomic information content \\
Winding conservation & $\leftrightarrow$ & genome integrity under attack \\
OHC defect self-assembly & $\leftrightarrow$ & capsid self-assembly \\
Condensate redundancy & $\leftrightarrow$ & capsid surface epitope display \\
\bottomrule
\end{tabular}
\end{center}

\textit{Proof sketch.} The icosahedral geometry of the capsid is
the 3D projection of D4 (Section~2). The D4 geometry is the BCT
vacuum structure. The capsid self-assembles from protein subunits
whose van der Waals contact geometry is BCT-constrained (Letter~67).
The genome nucleates capsid assembly exactly as topological charge
nucleates Hopf defect formation in the OHC. The closed shell
prevents genome escape exactly as the Hopf bounding sphere prevents
topological charge dispersal. The self-assembly free energy minimum
corresponds to the D4-symmetric (icosahedral) configuration because
D4 is the densest 4D lattice packing (BCT Uniqueness Theorem) and
its 3D projection maximises the enclosed volume for a given shell
area --- precisely the condition for capsid stability. $\square$

\section{Geometric Necessity vs.\ Convergent Evolution}

The standard view holds that icosahedral symmetry represents
convergent evolution: many independent viral lineages arrived at
icosahedral capsids because it is an evolutionary optimum~\cite{crick1956}.

The BCT-Virus Theorem establishes a stronger claim: icosahedral
symmetry is \textit{geometrically necessary}. Given:
\begin{enumerate}
\item Protein subunits whose contact geometry is BCT-constrained
  (van der Waals radii proportional to $a_0$, Letter~67).
\item The requirement to form a closed, stable, maximally symmetric
  shell around a genome.
\item The constraint that the shell be constructible from identical
  or quasi-identical subunits (genetic economy).
\end{enumerate}
The icosahedron is the \textit{unique} solution. No other regular
polyhedron satisfies all three constraints simultaneously. The
tetrahedron, cube, and octahedron have too few faces for large
genomes. The dodecahedron cannot tile a surface with triangular
subunits. Only the icosahedron, built from 20 equilateral triangular
faces, satisfies the geometric constraints imposed by D4 lattice
projection and BCT van der Waals contact geometry.

Evolution did not discover icosahedral geometry. It had no choice
but to use it. The universe has exactly one geometric template for
stable, information-containing, protein-assembled containers of this
topology --- and that template is D4.

\section{The ZV Geometric Container Conjecture}

We record the following conjecture, which extends the BCT-Virus
Theorem to all biological containers:

\textit{ZV Geometric Container Conjecture (March 2026).}
The universe employs a single geometric template --- D4 root lattice
projection --- for all stable, information-containing containers
across all length scales. This template manifests as:
\begin{itemize}
\item OHC Hopf defects at the Planck/QCD scale
  ($\sim 10^{-35}$--$10^{-15}$~m): topological charge enclosed by
  Hopf fibration geometry.
\item Viral capsids at the nanoscale ($\sim 10^{-8}$~m): genomic
  information enclosed by icosahedral D4-projection geometry.
\item Cellular organelles at the microscale ($\sim 10^{-6}$~m):
  biochemical information compartmentalised by membrane geometries
  selected by D4-constrained protein assembly.
\item (Conjectured) Cells themselves as the next scale.
\end{itemize}

\textit{Supporting evidence for organelles:}

\textbf{Clathrin-coated vesicles.} Clathrin assembles into icosahedral
cages --- the same T=1 and T=3 Caspar-Klug structures as small
viruses --- to bud vesicles from cellular membranes. The clathrin
triskelion (three-legged protein) assembles into the same 60-subunit
D4-projected geometry as viral capsids. Clathrin vesicles are
cellular information-routing containers built from D4 geometry.

\textbf{ATP synthase.} The F1 subunit of ATP synthase has C3
(three-fold) rotational symmetry; the Fo subunit has C10 symmetry
in mitochondria. The C3$\times$C10 interaction geometry is a
D4-related structure: $3 \times 10 = 30$, and 30 is the number of
2-fold symmetry axes of the icosahedron. The rotary mechanism of
ATP synthesis --- the fundamental energy currency of all life ---
exploits D4-projected symmetry mismatch to drive proton-coupled
rotation.

\textbf{Nuclear pore complexes.} The nuclear pore complex has C8
(eight-fold) symmetry: 8 outer rings, 8 inner rings, 8 spokes.
$8 = 2^3$ is a D4-related symmetry (the D4 lattice has 8 nearest
neighbours in its 4D structure). The nuclear pore --- the gateway
controlling all information flow between nucleus and cytoplasm ---
is D4-symmetric.

These three organelle examples suggest the ZV Conjecture is correct:
D4 geometry is life's universal template for containers that handle
information.

\section{Predictions}

\textbf{Prediction \#127}: The T-number series of viral capsids is
exactly the Eisenstein norm series $T = h^2 + k^2 + hk$. No stable
icosahedral capsid with non-Eisenstein T-value exists. This is
already confirmed for all known capsids.

\textbf{Prediction \#128}: The contact distance between capsid
protein subunits at pentameric vertices equals $5a_0(1+\alpha_0)
= 3.40$~\AA\ (Prediction~\#123), and at hexameric faces equals
$6a_0(1+\alpha_0) = 4.08$~\AA. These are BCT-predicted contact
distances testable by cryo-EM structural analysis of capsid protein
interfaces.

\textbf{Prediction \#129}: Clathrin cage assembly produces
preferentially T=1 ($60$ triskelions) and T=3 ($180$ triskelions)
structures --- identical to the smallest viral capsids --- because
both are governed by identical D4 projection geometry. No stable
clathrin cage with non-CK T-number exists.

\textbf{Prediction \#130}: The ZV Conjecture is falsified if any
stable, naturally occurring, protein-assembled information container
is found with non-icosahedral, non-D4-projected symmetry as its
preferred free-energy minimum. (Helical viruses are not
counterexamples: helical assembly is the 1D projection of D4,
preserving the translational symmetry of the root lattice.)

\section{Summary}

Viruses are macroscopic Hopf defects. The icosahedral geometry of
viral capsids is the 3D projection of the D4 root lattice, the same
structure that BCT derives as the unique quantum vacuum geometry.
The Caspar-Klug T-numbers are the norms of Eisenstein integers ---
the 2D sublattice norms of D4. The capsid self-assembly is the
macroscopic analogue of Hopf defect nucleation by topological charge
in the OHC. Icosahedral symmetry is not convergent evolution. It is
geometric necessity. The universe has one template for stable
information-containing containers, from Planck-scale quantum defects
to nanoscale viral capsids to microscale organelles. That template
is D4.

\begin{thebibliography}{99}
\bibitem{caspar1962} D.~L.~D.~Caspar, A.~Klug, ``Physical principles
  in the construction of regular viruses,''
  \textit{Cold Spring Harbor Symp.\ Quant.\ Biol.}\ \textbf{27},
  1 (1962).
\bibitem{crick1956} F.~H.~C.~Crick, J.~D.~Watson, ``Structure of
  small viruses,'' \textit{Nature}\ \textbf{177}, 473 (1956).
\bibitem{prl} M.~R.~Cabri\'e, BCT PRL Letter (2026), Zenodo:18884415.
\bibitem{monograph} M.~R.~Cabri\'e, BCT Monograph (2026),
  Zenodo:18885133.
\bibitem{bctjh} M.~R.~Cabri\'e, BCT Appendix~JH (2026),
  Zenodo:18975018.
\bibitem{bct64} M.~R.~Cabri\'e, BCT Letter~64 (2026)---OHC as
  Origin of Reality.
\bibitem{bct65} M.~R.~Cabri\'e, BCT Letter~65 (2026)---Geometry
  of Suffering.
\bibitem{bct66} M.~R.~Cabri\'e, BCT Letter~66 (2026)---Geometry
  of Life.
\bibitem{bct67} M.~R.~Cabri\'e, BCT Letter~67 (2026)---BCT-BIND.
\bibitem{zv2026} ZV, personal communication, March 2026.
  ``This isn't convergent evolution. This is geometric necessity.''
\end{thebibliography}

\end{document}
